% Section 4 — Related Work
% Status: v2 draft. Restructured per R3 venue notes — front-load
% the workshop audience's home turf (egocentric memory-augmented
% video LLMs), keep ANN/spatial-DB short. Citations curated against
% the workshop's CFP topic areas.

\section{Related Work}
\label{sec:related}

\paragraph{Egocentric video understanding and episodic memory.}
Ego4D~\cite{grauman2022ego4d} established look-back QA at scale and
EgoVLP~\cite{lin2022egovlp} the dominant egocentric video-language pretraining
recipe; EgoSchema~\cite{mangalam2023egoschema} exposed the temporal-grounding
collapse our prefilter targets. The ``where did I leave my keys?'' task
originates with Barmann et al.~\cite{barmann2022keys};
Di and Xie~\cite{di2024grounded} unified grounded egocentric QA and Chen et
al.~\cite{chen2025groundedmultihop} extended it to multi-hop chains.
Spatial cognition is central to cross-time localization~\cite{plizzari2025spatial};
EgoMask~\cite{liang2025egomask} benchmarks spatiotemporal grounding, and
ESOM's~\cite{manigrasso2026esom} 4\% online success rate underscores the
hardness of streaming egocentric retrieval. A contrasting spatial paradigm is
the learned world model: Navigation World Models~\cite{bar2025nwm} train an
implicit autoregressive generative model of egocentric space, predicting future
observations from past frames and actions. \psm{} is the explicit,
non-learned dual: an $O(\text{area})$-bounded spatial \emph{index} rather than
an unbounded learned predictor, and is orthogonal, a
bounded-memory spatial-temporal substrate any grounding method can rerank
without retraining (the VideoRAG pattern~\cite{jeong2025videorag}).

\paragraph{Memory-augmented video LLMs.}
MovieChat~\cite{song2024moviechat}, MA-LMM~\cite{he2024malmm},
Flash-VStream~\cite{zhang2024flashvstream},
LifelongMemory~\cite{wang2024lifelongmemory},
HERMES~\cite{zhang2026hermes}, InfiniPot-V~\cite{kim2025infinipotv}, and
SAVEMem~\cite{wu2026savemem} bound memory in \emph{token} space (KV-cache
compression, dense-to-sparse merging, recurrent state). \psm{} bounds
\emph{geographic} space (H3 cells) and runs upstream, deciding
\emph{which} frames the MLLM sees, not compressing what it sees.
Time-sensitive video LLMs such as TimeChat~\cite{ren2024timechat} instead target
temporal grounding within a fixed context, but do not bound memory, orthogonal
to \psm{}'s substrate role.

\paragraph{Streaming search and spatial indexing.}
DiskANN~\cite{jayaram2019diskann}, SPANN~\cite{chen2021spann}, and
locally-adaptive vector quantization~\cite{aguerrebere2024lvq} give
bounded-cost billion-point ANN but assume near-static corpora with no spatial
axis. \psm{} adds a ``near here'' H3 $k$-ring restriction and streaming
ring-buffer-aged per-cell state. Binning observations into a geohash or quadtree
grid and capping per bin already yields $O(\text{area})$ state; \psm{}'s
contribution is not the binning but what each cell carries: a time-decayed
HyperLogLog ring tracking cross-revisit cardinality and temporal envelopes (which
static buckets lack), reservoir-under-\caponek{} exemplar semantics, and streaming
age-out, all on H3's~\cite{brodsky2018h3} device-independent global ids.

\paragraph{Explicit spatial memory: maps and scene graphs.}
A parallel line builds \emph{explicit} spatial memories: open-vocabulary
semantic maps (VLMaps~\cite{huang2023vlmaps},
CLIP-Fields~\cite{shafiullah2023clipfields}) and online 3D scene graphs
(ConceptGraphs~\cite{gu2024conceptgraphs}, Hydra~\cite{hughes2022hydra}) attach
language-aligned features to reconstructed geometry. These are far richer scene
representations than \psm{}, but they assume per-scene SLAM/depth and grow with
scene content. \psm{} trades that richness for three deployment properties they
do not target: state \emph{provably} bounded by explored area, not scene
complexity; \emph{device-independent} localization via global H3 ids, with no
per-scene map to build or share; and a \emph{sensor-light} footprint (a coarse
pose stream, no dense reconstruction). It is a substrate such a builder could
sit atop, not a competitor.

\paragraph{Visual place recognition and loop closure.}
Recognizing a revisited place from appearance is the domain of visual place
recognition (VPR) and loop closure: classical bag-of-words
(FAB-MAP~\cite{cummins2008fabmap}, DBoW2~\cite{galvezlopez2012dbow2}) and learned
descriptors (CosPlace~\cite{berton2022cosplace},
EigenPlaces~\cite{berton2023eigenplaces}, MixVPR~\cite{alibey2023mixvpr},
AnyLoc~\cite{keetha2023anyloc}), purpose-built for the cross-session same-place
retrieval our HDD analysis (\S\ref{sec:results-hdd}) measures. We report that
$0.853$ AUC as a byproduct of the bounded index under off-the-shelf CLIP cosine,
not a VPR contribution: those descriptors and a $k$-nearest-cell control
(\S\ref{sec:limitations}) are the proper baselines, left to future work. \psm{}
differs in \emph{role}: it recognizes places to budget memory and route
retrieval, keyed on global H3 ids, so VPR descriptors are complementary drop-ins
for its cell-matching step, not alternatives.
